\scanpage{12}

Considérons la relation dans $I$
\[
i\unsure{<|\le}j\overset{\mathrm{déf}}{\Longleftrightarrow}\overline X_i\subset\overline X_j
\]
(Relation d'incidence des feuilles). C'est une relation de préordre, et grâce à 2) une relation d'ordre. Soit $\deleted{X_i}\inserted{X_{i_0}}$ une feuille minimale. Alors, comme $\overline X_{i_0}$ donc $\overline X_{i_0}-X_{i_0}$ est une réunion de feuilles, on voit que $\overline X_{i_0}\unsure{-}X_{i_0}=\varnothing$ à cause de la condition minimale : Une feuille minimale (pour la relation d'incidence) est fermée. Ceci étant, on raisonne par récurrence sur le nb des feuilles, en remplaçant $X$ par $X-X_{i_0}$ et la partition induite $(X_i)_{i\in I-\{i_0\}}$.

\margin{\uline{NB} Si $X_i$ n'est pas \uline{incidente} à $X_j$ alors $X_i\cap\overline X_j=\varnothing$, \deleted{donc} $\overline X_j\cap\overline X_i$ \uline{rare} dans $\overline X_i$.}

\inserted{Il est clair maintenant que $F_{\mathcal P}$ est engendré par ses parties fermées $F_0$, donc d) est satisfait automatiquement !}

\deleted{\unsure{Partition cellulaire de}} \inserted{\uline{Décomposition cellulaire}} (finie) de $X$ : satisfaisant $1^\circ)$ et $2^\circ)$.

Elle est ``\uline{stricte}'' si elle satisfait de plus

3) Les $X_i$ sont \deleted{\unsure{localement}}\inserted{connexes}.

Je dis que [moyennant $1^\circ)$ et $2^\circ)$] 3) équivaut à la condition f) pour $F_{\mathcal P}$, i.e. que $\unsure{\forall}$ $A\in F_{\mathcal P}$ qui est loc fermé est \deleted{\unsure{localement connexe|réunion loc connexe}} \inserted{a} des comp. connexes $\in F_{\mathcal P}$ \inserted{et en nb fini} \deleted{\uline{\unsure{C'est suffisant}}}. En effet,

$d)\Longrightarrow3)$, car $X_i$ étant loc$^t$ fermé et $\in F_{\mathcal P}$, \deleted{\unsure{alors}}

\mspage{13}
